\chapter{来源与史料说明}
\label{app:sources}

\section{来源索引的使用}
\label{app:sources:guide}

本附录按材料性质和时段提供入口，而不把不同文类排成固定的可靠性等级。事件索引中的链接回到正文首次展开处；图件索引只收录已经置入正文并有图注说明的示意图。遇到年代、数字、动机或制度范围的分歧，应回到相应事件的史料说明交叉阅读。

\section{基础叙事与编年材料}
\label{app:sources:narrative}

\begin{description}
  \sourceindexentry{汉书}{西汉纪年、人物、制度与地理问题的基础入口；与出土文献、后出的编年叙事和现代考证交叉使用。}
  \sourceindexentry{后汉书}{东汉政局、人物与制度材料的主要入口；应留意范晔编纂时的取舍及其与其他材料的距离。}
  \sourceindexentry{三国志}{魏、蜀、吴政治与人物叙事的基础材料；裴松之注保存的异文须分别辨认出处与成书背景。}
  \sourceindexentry{晋书}{西晋、东晋与十六国相关叙事的重要入口；不将后出修史的单一说法当作现场记录。}
  \sourceindexentry{宋书}{\sourcebook{南齐书}、\sourcebook{梁书}、\sourcebook{陈书}与\sourcebook{宋书}共同构成南朝政治、制度与人物材料的主要入口。}
  \sourceindexentry{魏书}{\sourcebook{北齐书}、\sourcebook{周书}、\sourcebook{北史}、\sourcebook{南史}与\sourcebook{魏书}共同构成北朝及南北互动的主要正史入口；各书的编纂立场与取材范围须随问题辨认。}
  \sourceindexentry{资治通鉴}{维持跨政权年序、比照前后关联的编年参照，不替代更接近具体问题的正史、文书或物质材料。}
\end{description}

\section{文书、宗教与物质材料}
\label{app:sources:documents}

\begin{description}
  \item[简牍与文书] 用于官府运作、律令、赋役、户籍、军政与地方实践等具体问题；出土地点、年代和保存状态决定其可说明的范围。
  \item[碑刻、墓志与造像题记] 可观察官职、家族、地方网络、纪念与宗教活动；不据个案直接推断全国通行的制度。
  \item[佛教、道教文献] 是理解信仰组织、仪式、文本流传和地方社会的重要入口，也须区分成书、抄写与后世重编的层次。
  \item[城址、墓葬、钱币与其他考古材料] 用于检验空间、生产、流通、聚落与物质文化条件；示意图不得把有限材料绘作精确疆界或完整路线。
\end{description}

\section{按时段进入材料}
\label{app:sources:periods}

\begin{description}
  \item[西汉与新莽] 从\sourcebook{汉书}进入纪年与制度线索，并以简牍、钱币、城址及相关现代研究核对行政与社会问题。
  \item[东汉] 从\sourcebook{后汉书}进入朝廷、地方与边疆叙事，并参照碑刻、简牍、经学和宗教材料。
  \item[三国与西晋] 从\sourcebook{三国志}、\sourcebook{晋书}和\sourcebook{资治通鉴}比照政权转换、战争与制度延续。
  \item[东晋、十六国与南北朝] 以晋书、南北朝正史和\sourcebook{资治通鉴}维持年序，并与墓志、造像、佛道文献和考古材料交叉阅读。
\end{description}

\section{与事件、图件索引的连接}
\label{app:sources:links}

可由\hyperref[app:index:events]{“事件索引”}回到正文锚点，由\hyperref[app:index:figures]{“图件索引”}查看相关示意图。新增事件条目使用 \verb|\eventindexentry{内部ID}{中文名称}|，新增图件条目使用 \verb|\figureindexentry{图件标签}{图名}|；只有目标锚点或标签已经写入正文时，才加入索引。
